% ===================================================================
% Section 2: Portfolio problem and design-based regularization
% ===================================================================
\section{Portfolio problem and design-based regularization}
\label{sec:setup}

\subsection{The portfolio problem}
Let there be $N$ assets with random return vector $r\in\R^{N}$ and covariance
$\Sig=\operatorname{Cov}(r)\in\Sym^{N}$. The long-only, budget-constrained global
minimum-variance problem is
\begin{equation}
  \min_{\w\in\R^{N}}\ \w^\top\Sig\,\w
  \qquad\text{s.t.}\qquad \one^\top\w=1,\quad \w\ge0,
  \label{eq:gmv}
\end{equation}
optionally with per-asset box bounds $\ell_i\le w_i\le u_i$ and group (sector)
caps $\sum_{i\in g}w_i\le c_g$ for a collection of index sets $g$. Throughout we
work with an estimate $\Sighat$ supplied by any covariance estimator (sample,
Ledoit--Wolf, OAS, or exponentially weighted). We write
\begin{equation}
  \Wset=\Big\{\w\in\R^{N}\ :\ \one^\top\w=1,\ \ \ell\le\w\le u,\ \
  \textstyle\sum_{i\in g}w_i\le c_g\ \ \forall g\Big\}
  \label{eq:feasible}
\end{equation}
for the feasible set, which is a (possibly unbounded) convex polyhedron, and is
compact whenever the box bounds are finite---in particular in the long-only case
$\ell=0$, where $\Wset$ is contained in the unit simplex.

\subsection{From a design measure to portfolio weights}
In optimal design theory, a design is a probability measure
$p=(p_1,\dots,p_J)$ over a set of \emph{design points}, that is, regressor vectors
$\{\vv_1,\dots,\vv_J\}$, and the Fisher information matrix of the associated
linear regression model is
\begin{equation}
  M(p)=\sum_{j=1}^{J}p_j\,\vv_j\vv_j^\top .
  \label{eq:designinfo}
\end{equation}
The decision vector $p$ lives on the probability simplex
$\{p\ge0,\ \one^\top p=1\}$, which is exactly the long-only budget set of
\eqref{eq:gmv}. This is the first half of the bridge: \emph{portfolio weights are
a design measure.} Substituting $\w$ for $p$ and attaching a vector $\vv_i$ to
each asset gives the \textbf{weight-dependent information matrix}
\begin{equation}
  \M \;=\; \sum_{i=1}^{N} w_i\,\vv_i\vv_i^\top \;=\; V\diag(\w)V^\top
  \;\in\;\Sym^{K},
  \qquad V=[\vv_1,\dots,\vv_N]\in\R^{K\times N},
  \label{eq:infomat}
\end{equation}
which is \emph{linear} in $\w$ and positive semidefinite whenever $\w\ge0$.

\begin{remark}[The one genuine disanalogy]
$\M$ is linear in $\w$, whereas portfolio variance $\w^\top\Sig\w$ is quadratic
in $\w$: the design objects are linear because each observation contributes an
independent rank-one term, while portfolio risk carries cross terms
$w_iw_j\Sig_{ij}$. The bridge is therefore exact at the level of \emph{matrices
and criteria} and heuristic at the level of the weight vector. Accordingly we use
the design criteria as \emph{regularizers added to} the quadratic risk objective,
not as a replacement for it.
\end{remark}

\subsection{What is $\vv_i$? A factor-regression interpretation}
\label{sec:vi}
The object \eqref{eq:infomat} acquires statistical meaning only once the
underlying regression is specified. We adopt a factor model for returns,
\begin{equation}
  r_t \;=\; B f_t + \varepsilon_t,
  \qquad B\in\R^{N\times K},\quad
  \E[f_t]=\mu_f,\quad \operatorname{Cov}(\varepsilon_t)=D,
  \label{eq:factor}
\end{equation}
where the $i$-th row $\bb_i^\top$ of the loading matrix $B$ collects asset $i$'s
exposures to $K$ common factors, $f_t$ is the vector of factor returns, and
$\varepsilon_t$ is idiosyncratic noise; this is the workhorse structure of
empirical asset pricing \cite{fama1993}.

Taking the design points to be the \textbf{factor loadings},
$\vv_i=\bb_i$---equivalently $V=B^\top$---the target parameter is
$\theta=\mu_f$, the vector of \textbf{factor premia}, and
\begin{equation}
  \M \;=\; \sum_{i=1}^{N} w_i\,\bb_i\bb_i^\top \;=\; B^\top\diag(\w)B .
  \label{eq:factorinfo}
\end{equation}
Interpretation: $\M$ is, up to scale, the Fisher information that a portfolio
holding $\w$ carries about the factor premia, and $\M^{-1}$ is the dispersion of
the corresponding estimates. The A/D/E criteria of $\M$ therefore control
\emph{how precisely, and how evenly, the portfolio identifies the common
factors}. The design points are genuine regressors, not a formal device.

\begin{example}[Two canonical choices of $V$]
\label{ex:Vchoices}
\begin{enumerate}[leftmargin=1.6em,itemsep=0.2em]
  \item \emph{Explicit factors, $K<N$.} $V=B^\top$ with loadings from an
    observable factor model (for instance Fama--French) or from the leading $K$
    principal components of $\Sighat$. Here $\M$ is $K\times K$ and genuinely
    rank-deficient relative to the number of assets.
  \item \emph{Precision metric (implicit factors), $K=N$.}
    $V=\Sighat^{-1/2}$, normalized as in Assumption~\ref{as:norm} below, so that
    $\M=\Sighat^{-1/2}\diag(\w)\Sighat^{-1/2}$ is $N\times N$ and nonsingular for
    $\w>0$. This is the limiting case in which the ``factors'' are the principal
    directions of the risk model itself; it requires no data beyond $\Sighat$ and
    is the specification used in the experiments of
    Sections~\ref{sec:computation}--\ref{sec:applications}.
\end{enumerate}
\end{example}

These two cases behave differently, and the difference is not cosmetic: several
familiar simplifications of the design criteria are valid only in the second. We
therefore develop the dependence on $V$ before defining the penalties.

\subsection{How the criteria depend on the choice of $V$}
\label{sec:Vchoice}

Write the three criteria as functions of the information matrix,
\begin{equation}
  \phiA(M)=\tr\!\big(M^{-1}\big),\qquad
  \phiD(M)=-\log\det M,\qquad
  \phiE(M)=-\lammin(M),
  \label{eq:criteria}
\end{equation}
each to be \emph{minimized}. (The signs are chosen so that all three are convex
functions of $M$ on the positive definite cone; see
Proposition~\ref{prop:convex}.)

\subsubsection{Invariance under change of factor basis}

The factor decomposition \eqref{eq:factor} is not unique: for any invertible
$T\in\R^{K\times K}$, writing $B\mapsto BT^\top$ and $f_t\mapsto T^{-\top}f_t$
leaves returns unchanged. Statistically estimated factors (principal components,
factor analysis) are identified only up to such a transformation, so it matters
whether the criteria depend on it.

\begin{theorem}[Basis invariance]
\label{thm:invariance}
Let $T\in\R^{K\times K}$ be invertible and replace $V$ by $TV$, so that
$M\mapsto TMT^\top$. Then, for every $\w$ with $\M\succ0$,
\[
  \phiD(TMT^\top)=\phiD(M)-2\log\lvert\det T\rvert ,
\]
so the D-criterion changes by an additive constant independent of $\w$; in
particular the D-penalized problem has the \emph{same} solution set for every
invertible $T$. By contrast,
\[
  \phiA(TMT^\top)=\tr\!\big(T^{-\top}M^{-1}T^{-1}\big),
  \qquad
  \phiE(TMT^\top)=-\lammin\!\big(TMT^\top\big),
\]
which in general differ from $\phiA(M)$ and $\phiE(M)$ as functions of $\w$. If
$T=Q$ is orthogonal, all three criteria are invariant:
$\phiA(QMQ^\top)=\phiA(M)$, $\phiE(QMQ^\top)=\phiE(M)$, and
$\phiD(QMQ^\top)=\phiD(M)$.
\end{theorem}

\begin{proof}
For the D-criterion, $\det(TMT^\top)=\det(T)\det(M)\det(T^\top)=(\det T)^2\det M$,
so $-\log\det(TMT^\top)=-\log\det M-2\log\lvert\det T\rvert$. The added term does
not depend on $\w$, so it shifts the objective by a constant and leaves the
argmin unchanged. For the A-criterion,
$(TMT^\top)^{-1}=T^{-\top}M^{-1}T^{-1}$, whence
$\tr((TMT^\top)^{-1})=\tr(T^{-\top}M^{-1}T^{-1})=\tr(M^{-1}T^{-1}T^{-\top})$,
which equals $\tr(M^{-1})$ for all $M$ if and only if $T^{-1}T^{-\top}=I$, that
is, if and only if $T$ is orthogonal. For the E-criterion, if $Q$ is orthogonal
then $QMQ^\top$ is similar to $M$ and has the same spectrum, so
$\lammin$ is unchanged; for general $T$ the spectrum is not preserved. The
orthogonal case of the D-criterion follows from $(\det Q)^2=1$.
\end{proof}

\begin{remark}[Modelling consequence]
\label{rem:rotation}
Theorem~\ref{thm:invariance} says that the D-penalty is \emph{well posed} under
statistically estimated factors, which are identified only up to an invertible
(indeed, up to an orthogonal) rotation, while the A- and E-penalties are well
posed only once the factor basis is fixed---for instance by taking principal
components with a fixed ordering and sign convention, or by using observable
factors. This is a practical guide to which criterion to use with which factor
construction, and to our knowledge it has not been noted in the portfolio
literature.
\end{remark}

\subsubsection{Scaling and why $V$ must be normalized}

\begin{proposition}[Scaling]
\label{prop:scaling}
Let $c>0$ and replace $V$ by $cV$, so that $M\mapsto c^{2}M$. Then
\[
  \phiA(c^{2}M)=c^{-2}\phiA(M),\qquad
  \phiE(c^{2}M)=c^{-0}\,\big({-}c^{2}\lammin(M)\big)=c^{2}\phiE(M),\qquad
  \phiD(c^{2}M)=\phiD(M)-2K\log c .
\]
Consequently the D-penalized problem is invariant to the scale of $V$, whereas
rescaling $V$ is equivalent to reparametrizing the penalty strengths as
$\gA\mapsto c^{-2}\gA$ and $\gE\mapsto c^{2}\gE$.
\end{proposition}

\begin{proof}
Immediate from $(c^2M)^{-1}=c^{-2}M^{-1}$, $\lammin(c^2M)=c^2\lammin(M)$, and
$\det(c^2M)=c^{2K}\det M$.
\end{proof}

Proposition~\ref{prop:scaling} is the reason a normalization convention is not
optional: without one, the numerical value of $\gA$ or $\gE$ has no meaning
across universes of different size or different return scales, and the penalty
strengths are not comparable between the two specifications of
Example~\ref{ex:Vchoices}. We adopt Assumption~\ref{as:norm} below, which fixes
$c$ by requiring $\tfrac1N\tr(VV^\top)=\tfrac1N\sum_i\lVert\vv_i\rVert_2^2=1$, and
thereby makes $\gA$ and $\gE$ scale-free.

\subsubsection{The rank-deficient case $K<N$: a Cauchy--Binet representation}

When $K<N$ the information matrix $\M=B^\top\diag(\w)B$ is $K\times K$ while the
weight vector has $N$ entries, so the map $\w\mapsto\M$ is not injective and the
determinant does not factor into a product over assets. The correct expansion is
combinatorial.

For an index set $S\subseteq\{1,\dots,N\}$ with $\lvert S\rvert=K$, let
$V_S\in\R^{K\times K}$ denote the submatrix of $V$ formed by the columns indexed
by $S$ (equivalently, the rows of $B$ indexed by $S$).

\begin{theorem}[Cauchy--Binet representation of the D-criterion]
\label{thm:cauchybinet}
Let $V\in\R^{K\times N}$ with $K\le N$ and let $\w\ge0$. Then
\begin{equation}
  \det \M \;=\; \sum_{\substack{S\subseteq\{1,\dots,N\}\\ \lvert S\rvert=K}}
  \Big(\prod_{i\in S}w_i\Big)\,\big(\det V_S\big)^{2},
  \label{eq:cb}
\end{equation}
and hence
$\phiD(\M)=-\log\sum_{\lvert S\rvert=K}\big(\prod_{i\in S}w_i\big)(\det V_S)^2 .$
\end{theorem}

\begin{proof}
Write $A=V\diag(\w)^{1/2}\in\R^{K\times N}$, so that $AA^\top=V\diag(\w)V^\top=\M$.
The Cauchy--Binet formula gives
$\det(AA^\top)=\sum_{\lvert S\rvert=K}\det(A_S)^2$, where $A_S$ is the
$K\times K$ submatrix of $A$ with columns indexed by $S$. Since
$A_S=V_S\diag(\w_S)^{1/2}$ we have
$\det(A_S)=\det(V_S)\prod_{i\in S}w_i^{1/2}$, and squaring and summing yields
\eqref{eq:cb}.
\end{proof}

Two consequences follow directly, and both matter for the formulation.

\begin{corollary}[Support condition for nonsingularity]
\label{cor:spanning}
Let $\rank V=K$ and $\w\ge0$ with support
$\supp(\w)=\{i:w_i>0\}$. Then $\M\succ0$ if and only if $\supp(\w)$ contains an
index set $S$ with $\lvert S\rvert=K$ and $\det V_S\neq0$; equivalently, if and
only if $\{\vv_i:w_i>0\}$ spans $\R^{K}$. In particular at least $K$ assets must
carry strictly positive weight.
\end{corollary}

\begin{proof}
$\M=AA^\top\succeq0$ always, so $\M\succ0$ if and only if $\det\M>0$. Every term
in \eqref{eq:cb} is nonnegative, so the sum is strictly positive if and only if
at least one term is, that is, if and only if there exists $S\subseteq\supp(\w)$
with $\lvert S\rvert=K$ and $\det V_S\neq0$. The latter is precisely the
statement that the columns $\{\vv_i\}_{i\in\supp(\w)}$ contain a basis of
$\R^{K}$, which requires $\lvert\supp(\w)\rvert\ge K$.
\end{proof}

\begin{corollary}[The square case: barrier reductions]
\label{cor:KN}
Suppose $K=N$ and $V$ is nonsingular. Then the only admissible index set in
\eqref{eq:cb} is $S=\{1,\dots,N\}$, so
\begin{equation}
  \phiD(\M)=-\log\det\M=-2\log\lvert\det V\rvert-\sum_{i=1}^{N}\log w_i ,
  \label{eq:Dbarrier}
\end{equation}
and, writing $G=(V^\top V)^{-1}$,
\begin{equation}
  \phiA(\M)=\tr\!\big(\M^{-1}\big)
  =\tr\!\big(\diag(\w)^{-1}G\big)
  =\sum_{i=1}^{N}\frac{G_{ii}}{w_i}.
  \label{eq:Abarrier}
\end{equation}
That is, up to constants the D-criterion is the classical \emph{logarithmic
barrier} and the A-criterion is a weighted \emph{inverse barrier} on the weights.
\end{corollary}

\begin{proof}
The first claim is immediate from \eqref{eq:cb} with a single term
$\big(\prod_iw_i\big)(\det V)^2$. For the second, $V$ nonsingular gives
$\M^{-1}=V^{-\top}\diag(\w)^{-1}V^{-1}$, so by cyclicity of the trace
$\tr(\M^{-1})=\tr\big(\diag(\w)^{-1}V^{-1}V^{-\top}\big)
=\tr\big(\diag(\w)^{-1}(V^\top V)^{-1}\big)$, which is
$\sum_iG_{ii}/w_i$.
\end{proof}

\begin{remark}[Scope of the barrier reductions]
\label{rem:scope}
Corollary~\ref{cor:KN} is stated deliberately as a \emph{special case}. The
identifications ``D-optimality $=$ log-barrier'' and ``A-optimality $=$ inverse
barrier'', which are the forms most often quoted, require $V$ to be square and
nonsingular, i.e.\ the precision-metric specification of
Example~\ref{ex:Vchoices}(2). Under a genuine factor model with $K<N$ they are
false, and the D-criterion must be handled through \eqref{eq:cb}. Both forms
remain convex in $\w$ (Proposition~\ref{prop:convex}), so the structural theory
of Section~\ref{sec:theory} covers both; but the interpretation differs. Under
$K=N$ the D-penalty acts separably on the weights and pushes toward $1/N$; under
$K<N$ it acts on the $K$-subsets of assets whose loadings form a well-conditioned
basis, and so rewards spreading weight across assets that are \emph{spanning} in
factor space rather than merely numerous. We regard the latter as the more
faithful expression of diversification, and it is a direction the empirical work
of this paper does not yet exploit.
\end{remark}
